\documentclass[11pt,a4paper]{scrartcl}
\usepackage{iutparakou}
\begin{document}

\fichetitre{Rapport — Jour 3}{Automates d'arbres (pivot) \& hash-consing}
{Binôme : DOSSOU Persévérance \& AGBOGBA Silas \quad$\bullet$\quad Modules : \texttt{formlang/tree.py} — Apps : \texttt{morpho, shield, hashcons}}

\section*{A. Réponses théoriques}

\subsection*{Q3.1 (0,5) — déterminisme \& coût}
Dans \texttt{morpho\_automaton} et \texttt{shield\_automaton}, les règles sont
enregistrées dans un dictionnaire \texttt{delta} indexé par la clé
\texttt{(symbole, états\_enfants)}. Un dictionnaire Python associe au plus une
valeur à une clé donnée : il ne peut donc jamais exister deux règles
$f(q_1,\dots,q_n)\to q$ et $f(q_1,\dots,q_n)\to q'$ avec $q\neq q'$ pour la même
clé — les deux automates sont \textbf{déterministes} par construction.
\emph{Conséquence :} pour un terme $t$, \texttt{TreeAutomaton.run} calcule un
unique état par nœud (aucun choix, aucun retour en arrière) ; le parcours
post-ordre visite chaque nœud exactement une fois $\Rightarrow$ coût
$O(n)$ où $n$ est le nombre de nœuds de $t$.

\subsection*{Q3.2 (0,5) — la frontière ne suffit pas}
Deux arbres $T_1=\texttt{word}(\texttt{prefixes}(\texttt{prefix(na)},
\texttt{nil}),\ \texttt{rest}(\texttt{root(pend)},\texttt{nil}))$ (classe
\texttt{PREFIXED}, frontière \texttt{na pend}) et $T_2$ où le même préfixe
\texttt{na} serait rattaché comme second niveau d'un circonfixe (structure
imbriquée différente, même séquence de feuilles \texttt{na pend}) ont une
frontière identique mais des classes différentes (\texttt{PREFIXED} vs.
\texttt{CIRCUMFIXED}, cf. énoncé §4.2.3). Un AFD de mots ne lit que la
\emph{séquence} des feuilles, jamais leur regroupement en sous-arbres ; il ne
peut donc pas distinguer $T_1$ de $T_2$, alors que le BUTA lit l'arbre
\emph{en profondeur} (règles \texttt{prefixes}/\texttt{rest}) et les classe
différemment.

\subsection*{Q3.3 (0,5) — théorème du yield}
Un mot morphologique accepté est de la forme \texttt{word(pc, rest(root, sc))}
où \texttt{pc} est une chaîne \texttt{prefixes} de $p\geq 0$ préfixes (état
\texttt{PFCHAIN} ou \texttt{NIL}) et \texttt{sc} une chaîne \texttt{suffixes}
de $s\geq 0$ suffixes. La frontière (yield) est la concaténation, dans
l'ordre de lecture gauche-droite, des labels des feuilles \texttt{prefix},
\texttt{root}, \texttt{suffix} : exactement $p^{*}\,r\,s^{*}$ (un préfixe
répété, une racine, un suffixe répété) — c'est \emph{le même langage
régulier} $p^{*}rs^{*}$ que celui de l'AFD du Jour 1 (§2, avec $p,r,s$ jouant
le rôle de $a,o,r$). Le BUTA généralise donc bien l'automate fini : sa
frontière retombe sur le langage régulier déjà reconnu par le DFA/NFA du
Jour 1.

\subsection*{Q3.4 (0,5) — réduplication}
La réduplication $\{ww \mid w\in\Sigma^*\}$ n'est pas un langage d'arbres
régulier. Un BUTA ne dispose que d'un nombre fini d'états $Q$ ; pour accepter
$ww$, l'automate devrait, en lisant la seconde moitié, « comparer » chaque
lettre à la lettre correspondante de la première moitié déjà lue — or l'état
courant ne peut encoder qu'une information bornée (au plus $|Q|$
possibilités), indépendante de $|w|$. Pour $|w|$ suffisant, deux mots
distincts $w\neq w'$ de même longueur mèneraient au même état après lecture
de la première moitié, et l'automate ne pourrait plus distinguer $ww$ de
$ww'$ : un argument de pompage sur les états (analogue au lemme de l'étoile)
montre l'absence de reconnaissance. Culy (1985) illustre précisément ce
phénomène pour le bambara, dont le vocabulaire dépasse le hors-contexte à
cause d'un motif de réduplication.

\section*{B. Exercices de programmation}

\begin{description}[leftmargin=1.6em,style=nextline]
  \item[E3.1] \texttt{TreeAutomaton.run} : étiquetage post-ordre récursif
    (état \texttt{REJECT} propagé dès qu'un enfant rejette ou qu'aucune règle
    ne correspond) ; \texttt{accepts} teste l'appartenance de l'état racine à
    \texttt{self.final}. \emph{Tests \texttt{test\_morpho\_*}, etc. débloqués.}
  \item[E3.2] Règles $\Delta$ de \texttt{morpho\_automaton} : feuilles
    \texttt{prefix/root/suffix/nil}, chaînage \texttt{prefixes}/\texttt{suffixes},
    \texttt{rest} (racine + suffixes), \texttt{word} (préfixes + rest) $\to$
    \texttt{WORD}. \emph{Tests \texttt{test\_morpho\_accept\_et\_classes},
    \texttt{test\_morpho\_rejet} : passés.}
  \item[E3.3] \texttt{\_seq(x,y)} : \texttt{DANGER} si l'un des deux est déjà
    \texttt{DANGER} ou si $\{x,y\}=\{\texttt{OVR},\texttt{ROLE}\}$, sinon la
    sévérité maximale ; \texttt{frame}/\texttt{sys} font passer un
    \texttt{OVR} ou \texttt{ROLE} directement encapsulé à \texttt{DANGER}.
    \emph{Test \texttt{test\_shield\_blocage} : passé.}
  \item[E3.4] \texttt{product(a1,a2)} : états = paires $(q_1,q_2)$, une règle
    $(sym,(q_1,q_2)\dots)\to(r_1,r_2)$ pour chaque paire de règles de même
    arité sur le même symbole ; finaux = $F_1\times F_2$. \emph{Test
    \texttt{test\_produit\_intersection} : passé.}
  \item[E3.5] \texttt{CompactStore.intern} : internement récursif des
    enfants, clé canonique \texttt{(symbol, label, kids\_ids)}, partage via
    table ; \texttt{get} reconstruit l'arbre exact (round-trip) ;
    \texttt{compression} $=1-\text{uniques}/\text{total}$. \emph{Tests
    \texttt{test\_round\_trip\_exact}, \texttt{test\_sous\_arbres\_identiques\_partages},
    \texttt{test\_compression\_positive} : passés.}
  \item[E3.6 (bonus)] \texttt{enc\_automaton} : compte les feuilles
    \texttt{enc} plafonné à 2 (même signature que \texttt{shield\_automaton})
    ; \texttt{dangerous\_and\_double\_encoded} = \texttt{product(shield\_automaton(),
    enc\_automaton())}. \emph{Test \texttt{test\_shield\_double\_encodage\_P45} :
    passé.}
\end{description}

\begin{exemple}[Sortie de référence obtenue]
\ttfamily \$ pytest -q tests/test\_tree.py tests/test\_hashcons.py\\ 8 passed
\end{exemple}

\end{document}
